\documentclass{beamer}
\input{tbook-deck}
\title{Stirling minors and Schur polynomials}
\author[Ma; Zhang]{Haobo Ma\inst{1} \and Wenlin Zhang\inst{1,2}}
\institute{\inst{1}Chrono AI, Singapore\\\inst{2}National University of Singapore, Singapore}
\date{Mathesis version 2 · 5 November 2026}
\begin{document}
% Manuscript: opening story, source introduction
\begin{frame}
  \titlepage
  \note{Begin with $x(x+1)(x+2)=x^3+3x^2+2x$ and the coefficient $c(3,2)=3$. Kurkov's A396563 asks why an alternating nested sum reproduces a Stirling determinant.}
\end{frame}
% Manuscript: story.md, question
\begin{frame}{The question: why does the nested sum work?}
  \begin{block}{A familiar start}
    \centering
    $x(x+1)(x+2)=x^3+3x^2+2x$,\qquad $c(3,2)=3$.
  \end{block}
  \vspace{4pt}
  \begin{alertblock}{Research question}
    Put $c(n+m+i,n+j)$ into a square matrix:
    \[
      B(n,k,m)=\det\!\left[c(n+m+i,n+j)\right]_{1\le i,j\le k}.
    \]
    Why should its value equal a nested sum of signs, factorials, and squares of differences?
  \end{alertblock}
  \note{Describe the object as a determinant of a coefficient table before introducing specialist language. The case $m=2$ is OEIS A396563; we work with general $m$.}
\end{frame}
\section{The bridge}
\subsection{Determinant to Schur}
% Manuscript: statements.md, s1
\begin{frame}{The identity and its Schur interpretation}
  \begin{bookresult}{s1-literal}
    For $n,k\ge0$ and $m\ge1$,
    \[ B(n,k,m)=F(n,m;n+k), \]
    where $F$ is the normalized alternating chain sum.
  \end{bookresult}
  \begin{block}{Written Schur interpretation}
    \[ B(n,k,m)=s_{(k^m)}(1,\ldots,n+m). \]
  \end{block}
  \note{The Lean declaration proves only the literal determinant-to-chain equality for $m\ge1$. The Schur equality is prior work of Merca and Amdeberhan--Ekhad, proved in the manuscript; it has its own formalization request. Define $F$ by the precise reading before using its abbreviation.}
\end{frame}
% Manuscript: story.md, proof mechanism
\begin{frame}{Proof mechanism: a triangular basis change}
  \begin{enumerate}
    \item Each entry is a coefficient of one common polynomial times a monic lower-degree polynomial.
    \item Replacing those polynomials by $1,z,z^2,\ldots$ is a triangular change with determinant $1$.
    \item The new determinant contains elementary symmetric functions of $1,2,\ldots,N$.
  \end{enumerate}
  \begin{block}{Identification}
    Dual Jacobi--Trudi gives the rectangle $s_{(k^m)}$. Its semistandard tableaux explain positivity and integrality.
  \end{block}
  \note{Explain the object before naming it: a tableau is a legal filling of a box rectangle, weak along rows and strict down columns. Its entry product is summed over all fillings.}
\end{frame}
% Manuscript: story.md, Vandermonde mechanism
\begin{frame}{Why a Vandermonde square appears}
  \begin{block}{Subset expansion}
    The alternant chooses $m$ entries from $1,\ldots,N$; the denominator contains every chosen--unchosen difference.
    \[
      P(z)=\prod_{r=1}^{N}(z-r),\qquad P'(a)=(-1)^{N-a}(a-1)!(N-a)!.
    \]
  \end{block}
  \begin{block}{Square and chain}
    Pairing differences within the chosen set gives $\prod_{a<b}(b-a)^2$; set
    $a_t=n-j_t+t$ to obtain $n\ge j_1\ge\cdots\ge j_m\ge0$.
  \end{block}
  \note{The substitution is a bijection between increasing subsets and decreasing chains. It explains the signs, factorials, and square together, rather than by finite pattern matching.}
\end{frame}
\section{Corrections}
\subsection{Two witnesses}
% Manuscript: statements.md, s2 and s4
\begin{frame}{Two small witnesses expose the alphabet error}
  \begin{block}{Two written corrections}
    The subset weights require the alphabet size $N=a+b$. A distinct unnumbered correction uses $b+r$ in the second-kind Schur alphabet.\par\smallskip
    First: correct value $3$; zero-extended printed weights give $1$.\par
    Second: $S(4,2)=7$; the printed alphabet gives $25$.
  \end{block}
  \note{For the first witness, $(a,b,r)=(1,1,1)$ gives $-1+2\cdot2=3$, while zero-extending the printed weights gives $1$. For the second, $S(4,2)=7$ but the printed alphabet gives $25$. Both are written corrections, without Lean badges.}
\end{frame}
\section{Integer exponents}
\subsection{Reciprocal branch}
% Manuscript: statements.md, s3 and s6
\begin{frame}{Allow every integer exponent}
  \begin{block}{All integer exponents}
    For $n,m\ge1$, $N=n+m$, the value of $F(n,m;a)$ is below. In the last line $r\ge0$:
    \[
    \begin{aligned}
    a\ge n &: s_{((a-n)^m)}(1,\ldots,N),\\
    -m<a<n &: 0,\\
    a=-m-r &: \frac{(-1)^{nm}}{(N!)^{m+r}}s_{(r^n)}(1,\ldots,N).
    \end{aligned}
    \]
  \end{block}
  \note{Keep the boundary visible: when $n=0$, the value is $F(0,m;a)=(m!)^a$, not part of the vanishing interval. The zero regime comes from two equal rows in the alternant.}
\end{frame}
% Manuscript: statements.md, s6
\begin{frame}{Reciprocal substitution gives A132393}
  \begin{bookresult}{s6}
    Set $a=-m-1$ and invert every variable. The $nm$ cross differences contribute $(-1)^{nm}$, and complementary rectangles give
    \[
      \left[\begin{matrix}n+m+1\\m+1\end{matrix}\right]
      =(n+m)!\prod_{t=1}^{m}(n+t)^tR(n,m).
    \]
  \end{bookresult}
  \note{This is the second literal source identity, corresponding to the frozen KurkovNestedStirlingIdentity.result. It meets the negative branch of s3, but its evidence scope is recorded separately.}
\end{frame}
\section{Consequences}
\begin{frame}{A first-kind minor as second-kind moments}
  \begin{block}{Written Corollary 4}
    For $n,k\ge0$, $m\ge1$, $N=n+m$, define
    \[ A_{ij}=S(n+k+i+j+1,N)\qquad(0\le i,j<m). \]
    Then
    \[ B(n,k,m)=(-1)^{\binom m2}\det A. \]
  \end{block}
  \note{The moments $M_\ell=\sum_{t=1}^N t^\ell/P'(t)$ equal $S(\ell+1,N)$. Cauchy--Binet produces the squared Vandermonde expansion and the parent subset formula supplies its sign. This is a written result, distinct from the corrected alphabet.}
\end{frame}

\subsection{What follows}
% Manuscript: statements.md, s5 and s7
\begin{frame}{The same expansion yields generating functions}
  \begin{block}{Generating functions and companion identities}
    The subset expansion gives a rational row generating function and shift recurrence. Companion formulas include the necessary $1/k!$ factor, the row polynomial, two recurrences, and harmonic weighted sums.
  \end{block}
  \note{These are written consequences and attributed earlier identities. Exact rational checks cover only the finite ranges listed in checks.json; they do not replace an all-parameter proof.}
\end{frame}
% Manuscript: story.md, scope
\begin{frame}{Evidence boundary}
  \begin{block}{Frozen Lean evidence}
    The three-parameter A396563 nested identity;\\
    the reciprocal-power A132393 identity.
  \end{block}
  \begin{alertblock}{Written mathematics in this paper}
    The Schur evaluation, integer extension, moment determinant, corrections and companion results retain their own written proofs and contribution links.
  \end{alertblock}
  \note{End by separating written derivations from frozen declarations. The reading chapters carry the complete statements and contribution links.}
\end{frame}
\appendix
% Manuscript: verification ledger (generated)
\FormalVerification
\end{document}
